\section{Atributos del Sistema}
Este capítulo presenta los atributos de calidad que el sistema debe satisfacer, 
las llamadas \textit{-ilities}, las restricciones que limitan las decisiones de
diseño y desarrollo, y las alternativas tecnológicas que se evaluaron antes de
optar por el desarrollo a medida. Los atributos enumerados aquí se concretarán
posteriormente en requisitos no funcionales medibles en el capítulo 7.

\subsection{Atributos de Calidad}
\begin{itemize}
    \item \textbf{Confidencialidad $-$ Prioridad: Crítica.} 
    Atributo derivado de la seguridad pero específicamente referido a la
    protección frente a divulgación no autorizada de datos. Las ``notas privadas''
    de los acompañamientos profesionales sólo deben ser visibles para voluntarios con el rol y 
    nivel correspondiente.

    \item \textbf{Fiabilidad} 
    Capacidad del sistema para funcionar correctamente bajo 
    condiciones normales de uso, minimizando fallos y tiempos de inactividad. Se materializa en la
    integridad referencial de la base de datos, la atomicidad transaccional de
    las operaciones que afectan a múltiples entidades (por ejemplo, una entrega
    de ropa que decrementa el stock y crea un registro histórico) y la copia de
    seguridad periódica con prueba de restauración.

    \item \textbf{Trazabilidad} 
    La trazabilidad es una de las 
    motivaciones centrales del proyecto. Se materializa en el registro inmutable de creaciones,
    modificaciones y eliminaciones (con identificador del voluntario y marca de
    tiempo), en el histórico de cambios de las entidades sensibles.

    \item \textbf{Cumplimiento legal} 
    El sistema debe satisfacer las exigencias documentales y de trazabilidad que
    imponen dos marcos legales: Ley 1/2025.

    \item \textbf{Mantenibilidad.}  
    Se materializa en una arquitectura por capas con responsabilidades bien delimitadas, código
    estructurado en módulos cohesionados, documentación técnica actualizada,
    cobertura de pruebas suficiente para detectar regresiones.

    \item \textbf{Escalabilidad}
    El volumen de datos previsto es modesto: la coordinación estima un universo
    de menos de cien amigos atendidos simultáneamente, en torno a una veintena
    de voluntarios activos, y un volumen de transacciones diarias del orden de
    decenas (no de miles ni de millones). El sistema debe funcionar
    correctamente en este orden de magnitud sin necesidad de optimizaciones
    costosas, pero la arquitectura debe permitir un crecimiento moderado sin reescritura.

    \item \textbf{Disponibilidad}
    Una caída de varias horas durante la madrugada no impide la operativa. Sin embargo, debe
    estar disponible durante los turnos de actividad del centro. Con el plan adecuado de despliegue
    podemos asegurar una disponibilidad del 99.5\% durante las horas de actividad.
\end{itemize}


\subsection{Restricciones}
Las restricciones limitan las decisiones de diseño y desarrollo del sistema.

\subsubsection{Restricciones Tecnológicas}
\begin{itemize}
    \item \textbf{Backend en Java con Spring Boot:} La elección responde al conocimiento
    previo del desarrollador, a la madurez del ecosistema Spring para
    aplicaciones empresariales y a la disponibilidad de Spring Security como
    base sólida para la implementación de autenticación y autorización.

    \item \textbf{Frontend en Angular:} La elección responde a tres factores:
    \par
    \vspace{5mm}
    Primero,
    la idoneidad del framework
    para aplicaciones de formulario intensivas como la presente.
    \par
    \vspace{5mm}
    Segundo, a la sinergia probada con backend Spring
    Boot: Angular y Spring Boot constituyen una combinación full-stack
    ampliamente adoptada en entornos de producción, en la que el contrato
    entre ambos lados es una API REST sobre HTTPS y la seguridad se articula
    mediante Spring Security con tokens JWT validados en cada petición. Esta
    arquitectura permite que el frontend opere como una aplicación
    desacoplada (\textit{Single Page Application}) que consume servicios sin estado,
    lo que facilita el despliegue independiente de cada capa, la
    escalabilidad horizontal y la sustitución del frontend en el futuro sin
    reescritura del backend. 
    \par
    \vspace{5mm}
    Tercero, al ecosistema maduro de
    herramientas que Angular proporciona: enrutamiento, formularios
    reactivos con validación, interceptores HTTP (ideales para inyectar el
    token JWT en cada petición y para gestionar errores de autorización de
    forma centralizada), y soporte de pruebas unitarias y end-to-end
    integrado. La sinergia técnica entre Angular y Spring Boot está
    ampliamente documentada en la literatura de arquitectura full-stack
    contemporánea [10].
    % TODO:
    % Deberíamos detallar la sinergia con sequence diagrams en algún Apéndice de acuerdo a los flujos en:
    % https://medium.com/@sanikapatil0213/how-angular-and-spring-boot-form-a-production-ready-full-stack-architecture-1bd63f3a84a2

    \item \textbf{Base de datos PostgreSQL:} La elección responde a su integración nativa con Render además de 
    la disponibilidad de cifrado de columna nativo (pgcrypto) para proteger los datos sensibles, y su
    cumplimiento estricto del estándar SQL\@.

    \item \textbf{Empaquetado mediante Docker.} Garantiza reproducibilidad del entorno entre desarrollo, 
    pruebas y producción además de permitir una BD PostgreSQL local desechable idéntica al de producción sin instalar nada en el sistema.

    \item \textbf{Despliegue en Render.} Plataforma de alojamiento gestionado en la
    nube. Implica aceptar las restricciones de los planes contratados. Sin embargo, ofrece 
    un nivel de suscripción gratuito para pruebas o demos (30 días máximos). 

    \item \textbf{Comunicación frontend-backend exclusivamente mediante API REST sobre HTTPS.}
\end{itemize}


\subsubsection{Restricciones de Hardware y Entorno}
\begin{itemize}
    \item El sistema debe ser accesible desde dispositivos con navegadores web modernos, incluyendo 
    ordenadores de escritorio, portátiles, tabletas y smartphones. No se requiere compatibilidad con 
    navegadores obsoletos.
    \item El sistema debe operar en condiciones de conectividad doméstica
    estándar, es decir, no presupone del uso de un ancho de banda empresarial para su funcionamiento correcto.
\end{itemize}


\subsubsection{Restricciones Operativas}
\begin{itemize}
    \item \textbf{Auditoría obligatoria de operaciones sensibles.} Toda creación,
    modificación o eliminación de información relevante (especialmente datos
    sensibles, entregas, gastos, asientos contables y entradas alimentarias)
    debe quedar registrada con identificador del voluntario, marca de
    tiempo y, en caso de modificación, valor anterior.

    \item \textbf{Idioma único: español.} La interfaz, los mensajes de error, la
    documentación de usuario y los informes generados se producen en español de España.
\end{itemize}



\subsection{Alternativas Consideradas}
Antes de optar por el desarrollo a medida, durante las entrevistas se
evaluaron explícita o implícitamente varias alternativas. A continuación se
documentan las consideradas y los motivos que llevaron a descartarlas.


\subsubsection{Continuar con Hojas de Cálculo en Google Drive}
La alternativa más sencilla y económica habría sido sistematizar las hojas
de cálculo actuales en una estructura compartida en Google Drive. Esta
opción se descartó por \textbf{concurrencia limitada} (aunque \textit{Google Sheets} permite edición
simultánea, no soporta de forma fiable la operación concurrente de
múltiples voluntarios sobre las mismas filas), \textbf{carga manual elevada}, \textbf{dificultad para generar informes complejos} 
debido a la descentralización de la información en múltiples hojas, y \textbf{riesgo de pérdida de datos} por errores humanos e 
\textbf{imposibilidad de cumplir la Ley 1/2025} de forma cómoda y fiable.


\subsubsection{CiviCRM}
CiviCRM es un CRM de código abierto orientado a organizaciones sin ánimo
de lucro gratuito. Fue propuesto por uno de los
colaboradores de la directiva y evaluado durante la segunda entrevista más impone 
uso sistema operativo Linux que no se posee ni se desea. Además, gira en torno a
donantes, contribuciones y campañas de captación de fondos; su modelo
no se ajusta de forma natural al concepto de ``amigo con ficha
progresiva por nivel de confianza'' que es central al proyecto.


\subsubsection{Salesforce}
Salesforce ofrece una variante de su plataforma orientada a entidades sin
ánimo de lucro, con descuentos significativos y dotación gratuita inicial
de licencias (NPSP / Power of Us). Se evaluó y se descartó por los
siguientes motivos:
\begin{itemize}
    \item Aunque la dotación inicial es
    gratuita, es altamente limitada en cuanto al número de usuarios. 
    \item La ampliación de licencias, el almacenamiento adicional y el
    soporte de funcionalidades avanzadas implican costes recurrentes elevados.
    \item Ofrece una
    funcionalidad muy superior a la requerida, gran parte de la cual no
    se aprovecharía, lo que se traduce en una interfaz innecesariamente
    compleja para los voluntarios.
    \item Imposibilidad de unificar dominios de trazabilidad alimentaria
   y la contabilidad operativa para
    cuentas y mapeo de extractos bancarios; escapan al alcance natural de
    un CRM y exigirían integraciones con sistemas adicionales.
\end{itemize}


\subsubsection{ERPs de Código Abierto}
Se descartaron por motivos similares a los de CiviCRM y Salesforce: sobre-dimensionamiento,
inadecuación conceptual al dominio y dificultad para implementar el
módulo específico de trazabilidad alimentaria. Los consultados fueron ERPNext y Dolibarr.


\subsubsection{Conclusión}
La combinación de tres factores hace inviables las alternativas evaluadas y
justifica el desarrollo a medida:
\begin{enumerate}
    \item La necesidad de implementar el \textbf{módulo de trazabilidad alimentaria de la
    Ley 1/2025} específico del dominio de ayuda comunitaria en cualquier
    solución elegida.

    \item La necesidad de unificar en una misma aplicación la gestión de personas,
    servicios, inventario, alimentación y contabilidad, lo que ningún
    producto comercial ofrece de forma nativa.

    \item La restricción presupuestaria que excluye soluciones SaaS de coste
    mensual significativo.
\end{enumerate}